\section{Continuous Assessment 1 (CA-1) --- Set A}
\label{sec:ca1_set_a}

\LPUPaperHeader{Continuous Assessment (CA-1)}{A}{60 Minutes}{30 Marks}{CO1 (Matrix Algebra), CO2 (Homogeneous Differential Equations)}

\begin{center}
\sffamily\textbf{\color{AlertRed}Note: Attempt all six questions. Each question carries 5 marks.}
\end{center}

\questionitem{1}{
    Reduce the matrix $A$ to Row Echelon Form and determine its rank:
    \[
    A = \begin{bmatrix} 1 & 2 & 3 & 2 \\ 2 & 3 & 5 & 1 \\ 1 & 3 & 4 & 5 \end{bmatrix}
    \]
}{5}{CO1}{L3 Apply}

\questionitem{2}{
    Using Gauss-Jordan elimination, find the inverse of the matrix:
    \[
    A = \begin{bmatrix} 2 & 1 & 1 \\ 1 & 2 & 1 \\ 1 & 1 & 2 \end{bmatrix}
    \]
}{5}{CO1}{L3 Apply}

\questionitem{3}{
    Investigate the consistency of the system of equations for all real values of $\lambda$ and $\mu$:
    \[
    \begin{aligned}
    x + y + z &= 6 \\
    x + 2y + 3z &= 10 \\
    x + 2y + \lambda z &= \mu
    \end{aligned}
    \]
}{5}{CO1}{L4 Analyze}

\questionitem{4}{
    Find the eigenvalues and eigenvectors of the symmetric matrix:
    \[
    A = \begin{bmatrix} 3 & 1 & 1 \\ 1 & 3 & 1 \\ 1 & 1 & 3 \end{bmatrix}
    \]
}{5}{CO1}{L3 Apply}

\questionitem{5}{
    Verify the Cayley-Hamilton theorem for the matrix $A = \begin{bmatrix} 1 & 2 \\ 2 & -1 \end{bmatrix}$ and compute $A^4$.
}{5}{CO1}{L3 Apply}

\questionitem{6}{
    Solve the homogeneous differential equation: $(D^3 - 6D^2 + 11D - 6)y = 0$.
}{5}{CO2}{L3 Apply}

\newpage
\subsection*{Solutions and Marking Rubrics --- CA-1 Set A}

\begin{solutionbox}[Solution to Question 1]
Apply row operations: $R_2 \to R_2 - 2R_1$, $R_3 \to R_3 - R_1$:
\[
\begin{bmatrix} 1 & 2 & 3 & 2 \\ 0 & -1 & -1 & -3 \\ 0 & 1 & 1 & 3 \end{bmatrix} \xrightarrow{R_3 \to R_3 + R_2} \begin{bmatrix} 1 & 2 & 3 & 2 \\ 0 & -1 & -1 & -3 \\ 0 & 0 & 0 & 0 \end{bmatrix}
\]
The matrix has 2 non-zero rows in row echelon form. Therefore, $\mathbf{\operatorname{rank}(A) = 2}$.
\begin{markingrubric}
$\bullet$ Formulating row operations $R_2 - 2R_1$ and $R_3 - R_1$: \textbf{2 Marks}\\
$\bullet$ Reaching echelon form: \textbf{2 Marks}\\
$\bullet$ Concluding rank is 2: \textbf{1 Mark}
\end{markingrubric}
\end{solutionbox}

\begin{solutionbox}[Solution to Question 2]
Set up $[A \mid I]$ and row reduce to $[I \mid A^{-1}]$. Determinant is $\det(A) = 2(3) - 1(1) + 1(-1) = 4 \neq 0$.
Performing row reductions yields:
\[
A^{-1} = \frac{1}{4}\begin{bmatrix} 3 & -1 & -1 \\ -1 & 3 & -1 \\ -1 & -1 & 3 \end{bmatrix}
\]
\begin{markingrubric}
$\bullet$ Correctly setting up augmented matrix: \textbf{1 Mark}\\
$\bullet$ Pivot operations: \textbf{3 Marks}\\
$\bullet$ Final inverse matrix: \textbf{1 Mark}
\end{markingrubric}
\end{solutionbox}

\begin{solutionbox}[Solution to Question 3]
Augmented matrix:
\[
[A \mid B] = \begin{bmatrix} 1 & 1 & 1 & 6 \\ 1 & 2 & 3 & 10 \\ 1 & 2 & \lambda & \mu \end{bmatrix} \sim \begin{bmatrix} 1 & 1 & 1 & 6 \\ 0 & 1 & 2 & 4 \\ 0 & 0 & \lambda - 3 & \mu - 10 \end{bmatrix}
\]
\begin{enumerate}
    \item \textbf{No Solution}: $\lambda = 3$ and $\mu \neq 10$ ($\operatorname{rank}(A) = 2 \neq \operatorname{rank}([A\mid B]) = 3$).
    \item \textbf{Unique Solution}: $\lambda \neq 3$ ($\operatorname{rank}(A) = \operatorname{rank}([A\mid B]) = 3 = n$).
    \item \textbf{Infinitely Many Solutions}: $\lambda = 3$ and $\mu = 10$ ($\operatorname{rank}(A) = \operatorname{rank}([A\mid B]) = 2 < 3$).
\end{enumerate}
\begin{markingrubric}
$\bullet$ Echelon reduction: \textbf{2 Marks}\\
$\bullet$ Clear analysis of each of the three parametric cases: \textbf{3 Marks}
\end{markingrubric}
\end{solutionbox}

\begin{solutionbox}[Solution to Question 4]
Characteristic polynomial: $\lambda^3 - 9\lambda^2 + 24\lambda - 16 = 0 \implies (\lambda - 5)(\lambda - 2)^2 = 0$.
Eigenvalues: $\lambda_1 = 5$, $\lambda_2 = \lambda_3 = 2$.
Eigenvector for $\lambda = 5$: $X_1 = [1, 1, 1]^T$.
Eigenvectors for $\lambda = 2$: $X_2 = [-1, 1, 0]^T$, $X_3 = [-1, 0, 1]^T$.
\begin{markingrubric}
$\bullet$ Characteristic equation \& roots: \textbf{2 Marks}\\
$\bullet$ Eigenvector for $\lambda = 5$: \textbf{1.5 Marks}\\
$\bullet$ Independent eigenvectors for $\lambda = 2$: \textbf{1.5 Marks}
\end{markingrubric}
\end{solutionbox}

\begin{solutionbox}[Solution to Question 5]
Characteristic equation: $|A - \lambda I| = \lambda^2 - 5 = 0$.
By Cayley-Hamilton: $A^2 - 5I = 0 \implies A^2 = 5I$.
Therefore, $A^4 = (A^2)^2 = (5I)^2 = 25I = \begin{bmatrix} 25 & 0 \\ 0 & 25 \end{bmatrix}$.
\begin{markingrubric}
$\bullet$ Verifying Cayley-Hamilton identity: \textbf{2.5 Marks}\\
$\bullet$ Computing $A^4 = 25I$: \textbf{2.5 Marks}
\end{markingrubric}
\end{solutionbox}

\begin{solutionbox}[Solution to Question 6]
Auxiliary equation: $m^3 - 6m^2 + 11m - 6 = 0$.
Roots: $(m - 1)(m - 2)(m - 3) = 0 \implies m = 1, 2, 3$.
Since roots are real and distinct:
\[
\mathbf{y(x) = c_1 e^x + c_2 e^{2x} + c_3 e^{3x}}
\]
\begin{markingrubric}
$\bullet$ Writing auxiliary equation: \textbf{1.5 Marks}\\
$\bullet$ Factoring roots $1, 2, 3$: \textbf{2 Marks}\\
$\bullet$ General solution: \textbf{1.5 Marks}
\end{markingrubric}
\end{solutionbox}

\newpage
